\section{Discussion}

The main message of this study is not that neural networks dominate XGBoost on every metric. Rather, the evidence supports a more clinically meaningful interpretation: tactile spatiotemporal data can first support reliable intraoperative nodule detection, then support size inversion, and under explicit hierarchical organization, support coarse depth discrimination. This distinction matters because it aligns the modeling strategy with the actual priority order of intraoperative use.

The strength of the present work lies in the way the problem was organized. We did not force all targets into a flat multitask setting. Instead, we used the signal hierarchy revealed by mechanism analysis: detection as the strongest and most actionable signal, size as a strong secondary signal, and depth as a weaker conditional signal that must be modeled under size awareness. This task reorganization turned out to be more important than simply increasing network complexity.

The study also suggests that XGBoost and neural networks should not be framed as interchangeable competitors. XGBoost first answered whether the tactile data contained learnable information and which physical feature families were most relevant. The raw-input neural pipeline then answered whether the original tensors themselves could support learning without directly feeding those handcrafted physical features into the scientific main model. Finally, the unified hierarchical inverter addressed the engineering question of how to improve predicted-route robustness under realistic deployment conditions.

This role separation is important for interpretation. The raw-input hierarchical neural model should be regarded as the scientific main model because it provides the strongest evidence that original tactile tensors can support detection, size learning, and size-aware coarse depth discrimination, while also enabling interpretability analysis. By contrast, the unified hybrid inverter should be regarded as a deployment enhancement model. Because it explicitly integrates structured physical features, it is well suited for system robustness improvement, but it should not be used as the sole evidence that all physical structure was learned automatically from raw input.

Another practical insight is that the best intermediate metric is not equivalent to the best system-level route. For example, XGBoost still achieved the best standalone size regression MAE, but the unified hierarchical inverter produced the best predicted-route depth performance. This indicates that for intraoperative use, optimizing a single subtask in isolation is less important than improving the robustness of the full decision chain under real routing uncertainty.

Depth remains the most delicate part of the study. The evidence supports the existence of learnable depth information, but only at the level of coarse depth discrimination and only under explicit size-aware organization. The phase-occlusion results also suggest that the current raw-input model still relies more on peak-neighborhood information than on the full loading/release dynamics implied by the mechanism analysis. For this reason, we intentionally avoid framing the present work as a mature continuous depth inversion system.

The study has several limitations. First, the raw-input models did not surpass the structured baseline on every size metric, particularly in standalone regression MAE. Second, depth was limited to a three-class coarse formulation rather than continuous inversion. Third, the experiments were performed on ex vivo porcine lungs rather than clinical intraoperative data. Fourth, the interpretability results should be read as evidence of partial mechanistic consistency, not as a complete recovery of the underlying biomechanics. Future work should strengthen phase-aware modeling, improve route-noise robustness, and test the system under more clinically realistic variability.
